\documentclass[12pt,a4paper]{article}

\usepackage[utf8]{inputenc}
\usepackage{xepersian}
\settextfont{Yas}
\setmathdigitfont{Yas}

\usepackage{amsmath,amssymb,amsthm}
\usepackage{geometry}
\geometry{top=2.5cm,bottom=2.5cm,left=2.5cm,right=2.5cm}
\usepackage{booktabs}
\usepackage{array}
\usepackage{enumitem}
\usepackage{tcolorbox}
\usepackage{fancyhdr}
\usepackage{graphicx}
\usepackage{mdframed}

\pagestyle{fancy}
\fancyhf{}
\fancyhead[R]{آزمون میان‌ترم: هفته‌های ۱ تا ۴}
\fancyfoot[C]{\thepage}
\renewcommand{\headrulewidth}{0.4pt}

\newcommand{\bitem}[1]{\textbf{#1}}

\begin{document}

\begin{center}
{\Large\bfseries آزمون میان‌ترم: هفته‌های ۱ تا ۴}\\[4pt]
{\large یادگیری ماشین برای آمادگی المپیاد بین‌المللی هوش مصنوعی}\\[4pt]
زمان: ۹۰ دقیقه \quad|‌‌quad کتاب‌خانه بسته\\[2pt]
ماشین‌حساب مجاز است
\end{center}

\rule{\textwidth}{0.5pt}

\tableofcontents

\bigskip

\section*{بخش اول: سوالات تشریحی طولانی (۶۰ دقیقه)}

\subsection*{سوال ۱ [۲۰ نمره] — از چارچوب تا رگرسیون تنظیم‌شده}

داده‌های $\mathcal{D} = \{(x_1, y_1), (x_2, y_2), \ldots, (x_n, y_n)\}$ را با یک ویژگی ورودی $x$ و خروجی پیوسته $y$ در نظر بگیرید. از مدل خطی $\hat{y} = wx + b$ با تابع زیان خطای مربعی $L(\hat{y}, y) = (\hat{y} - y)^2$ استفاده می‌کنیم.

\bigskip

\textbf{(الف)} [۳ نمره] فرمول ریسک تجربی $R_{\text{emp}}(w, b)$ و ریسک واقعی $R(w, b)$ را بنویسید. اصل کمینه‌سازی ریسک تجربی (ERM) را در یک جمله بیان کنید. توضیح دهید چرا نمی‌توانیم ریسک واقعی را مستقیماً کمینه کنیم.

\bigskip

\textbf{(ب)} [۵ نمره] با ثابت فرض کردن $w$ و کامل کردن مربع در $b$، عرض از مبدأ بهینه $b^*$ را استخراج کنید. نشان دهید که $b^* = \bar{y} - w\bar{x}$ و در یک جمله توضیح دهید چرا این بدان معناست که خط بهینه همواره از نقطه میانگین‌ها $(\bar{x}, \bar{y})$ عبور می‌کند.

\bigskip

\textbf{(ج)} [۴ نمره] با جایگذاری $b^*$، تابع زیان به یک درجه دوم بر حسب $w$ تبدیل می‌شود:
$$R(w) = \text{Var}(x) \cdot w^2 - 2\text{Cov}(x, y) \cdot w + \text{Var}(y)$$
با استفاده از این، شیب OLS را به دست آورید: $w^* = \text{Cov}(x, y)/\text{Var}(x)$.

\bigskip

اکنون رگرسیون ریج را با تابع هدف زیر در نظر بگیرید:
$$R_{\text{ridge}}(w, b) = \frac{1}{n}\sum_{i=1}^{n}(wx_i + b - y_i)^2 + \lambda w^2$$

\textbf{(د)} [۴ نمره] جواب ریج $w^*_{\text{ridge}} = \frac{\text{Cov}(x,y)}{\text{Var}(x) + \lambda}$ است. $w^*_{\text{ridge}}$ را بر حسب $w^*_{\text{OLS}}$ بیان کنید و ضریب جمع‌شونده (shrinkage factor) را مشخص کنید. وقتی $\lambda \to 0$ و $\lambda \to \infty$ چه اتفاقی برای این ضریب می‌افتد؟

\bigskip

\textbf{(هـ)} [۴ نمره] جدول زیر را با استفاده از «پایین»، «متوسط» یا «بالا» پر کنید:

\begin{center}
\begin{tabular}{|l|c|c|c|}
\hline
 & $\lambda = 0$ (OLS) & $\lambda$ متوسط & $\lambda \to \infty$ \\
\hline
\textbf{بایاس} & ؟ & ؟ & ؟ \\
\hline
\textbf{واریانس} & ؟ & ؟ & ؟ \\
\hline
\textbf{خطای آموزش} & ؟ & ؟ & ؟ \\
\hline
\textbf{خطای تست} & ؟ & ؟ & ؟ \\
\hline
\end{tabular}
\end{center}

در ۲ تا ۳ جمله توضیح دهید: چرا رگرسیون ریج با وجود داشتن بایاس، گاهی خطای تست پایین‌تری نسبت به OLS دارد؟

\newpage

\subsection*{سوال ۲ [۲۰ نمره] — بیش‌برازش، انتخاب مدل و ارزیابی}

۲۰ نقطه داده از رابطه واقعی $y = 0.5x^2 + \epsilon$ که $\epsilon$ نوفه تصادفی است، تولید شده است. مدل‌های رگرسیون چندجمله‌ای با درجه‌های مختلف برازش داده شده‌اند:

\begin{center}
\begin{tabular}{|c|c|c|}
\hline
درجه $d$ & MSE آموزش & MSE تست \\
\hline
1 & 12.5 & 13.1 \\
\hline
2 & 1.8 & 2.3 \\
\hline
3 & 1.5 & 2.1 \\
\hline
5 & 0.6 & 4.7 \\
\hline
8 & 0.02 & 18.5 \\
\hline
15 & 0.001 & 32.0 \\
\hline
\end{tabular}
\end{center}

\textbf{(الف)} [۳ نمره] برای درجه‌های ۱، ۲ و ۱۵، مشخص کنید مدل کم‌برازش است، بیش‌برازش است یا برازش خوب است. هر پاسخ را با استفاده از مقادیر خطا و شکاف تعمیم توجیه کنید.

\bigskip

\textbf{(ب)} [۳ نمره] ثابت کنید که خطای آموزش تابعی نافزاینده از درجه چندجمله‌ای $d$ است. یعنی اگر $d_1 < d_2$ آنگاه $R_{\text{train}}(d_2) \leq R_{\text{train}}(d_1)$.\\
\textit{(راهنمایی: به فضاهای فرض $\mathcal{H}_{d_1}$ و $\mathcal{H}_{d_2}$ فکر کنید.)}

\bigskip

\textbf{(ج)} [۴ نمره] برای انتخاب بهترین درجه چندجمله‌ای از اعتبارسنجی متقاطع ۵-بخشی (5-fold CV) استفاده می‌کنید. ۲۰ نقطه داده دارید. در هر بخش آموزش چند نمونه هست؟ در هر بخش اعتبارسنجی چند نمونه؟ مجموعاً چند بار مدل برازش داده می‌شود؟ مراحل ۵گانه انتخاب مدل را توصیف کنید.

\bigskip

\textbf{(د)} [۴ نمره] فرض کنید به‌جای اعتبارسنجی متقاطع، ۵۰۰ درجه مختلف را امتحان می‌کنید و درجه‌ای با کمترین خطای اعتبارسنجی روی یک مجموعه اعتبارسنجی انتخاب می‌کنید. خطر «بیش‌برازش به مجموعه اعتبارسنجی» را توضیح دهید. چرا مجموعه تست (که فقط یک بار استفاده می‌شود) این مشکل را حل می‌کند؟

\bigskip

\textbf{(هـ)} [۳ نمره] داده‌های شما مؤلفه زمانی دارند (مثلاً مقادیر $x$ تاریخ هستند). قبل از تقسیم به آموزش/تست، داده‌ها را مخلوط (shuffle) می‌کنید. چه نوع نشت داده‌ای رخ می‌دهد و چرا مشکل‌ساز است؟ داده‌های سری زمانی را چگونه باید تقسیم کنید؟

\bigskip

\textbf{(و)} [۳ نمره] پس از انتخاب درجه $d=2$، روی مجموعه تست ارزیابی می‌کنید و $R^2 = 0.91$ به دست می‌آید. توضیح دهید $R^2 = 0.91$ چه معنایی دارد. اگر به‌جای آن $R^2 = -0.05$ می‌گرفتید، چه نشان می‌داد؟

\newpage

\subsection*{سوال ۳ [۲۰ نمره] — توابع زیان، هندسه تنظیم و معیارهای طبقه‌بندی}

\textbf{(الف)} [۴ نمره] دو تابع زیان برای رگرسیون را در نظر بگیرید: خطای مربعی $L_2(\hat{y}, y) = (\hat{y}-y)^2$ و خطای مطلق $L_1(\hat{y}, y) = |\hat{y}-y|$.

\begin{enumerate}[label=(\roman*)]
\item کدام یک نسبت به داده‌های پرت (outliers) حساس‌تر است؟ با یک مثال عددی توجیه کنید.
\item زیان هوبر (Huber) برای خطاهای کوچک به $L_2$ و برای خطاهای بزرگ به $L_1$ تبدیل می‌شود. در ۲-۳ جمله توضیح دهید چرا این یک مصالحه خوب است.
\end{enumerate}

\bigskip

\textbf{(ب)} [۵ نمره] تنظیم L1 (لاسو) و L2 (ریج) را در حالت دو پارامتر ($w_1, w_2$) در نظر بگیرید.

\begin{enumerate}[label=(\roman*)]
\item جمله جریمه هر یک را بنویسید.
\item ناحیه محدودیت هر یک را در صفحه $(w_1, w_2)$ توصیف کنید.
\item توضیح دهید چرا لاسو تمایل به ایجاد راه‌حل‌های تنک (برخی $w_j = 0$) دارد ولی ریج این‌گونه نیست. در پاسخ خود به هندسه ارجاع دهید.
\item یک سناریو که در آن لاسو بر ریج ترجیح داده می‌شود و یک سناریو که در آن ریج بر لاسو ترجیح داده می‌شود بیان کنید.
\end{enumerate}

\bigskip

\textbf{(ج)} [۵ نمره] یک فیلتر اسپم روی ۲۰۰۰ ایمیل (۲۰۰ اسپم، ۱۸۰۰ غیراسپم) ارزیابی می‌شود. ماتریس درهم‌ریختگی:

\begin{center}
\begin{tabular}{|l|c|c|}
\hline
 & پیش‌بینی اسپم & پیش‌بینی غیراسپم \\
\hline
\textbf{اسپم واقعی} & 160 & 40 \\
\hline
\textbf{غیراسپم واقعی} & 30 & 1770 \\
\hline
\end{tabular}
\end{center}

\begin{enumerate}[label=(\roman*)]
\item دقت (accuracy)، دقت مثبت (precision)، بازیابی (recall) و F1 را محاسبه کنید.
\item یک مدل پایه که برای همه چیز «غیراسپم» پیش‌بینی می‌کند چه دقتی دارد؟ توضیح دهید چرا این نشان می‌دهد دقت به‌تنهایی گمراه‌کننده است.
\item اگر این یک مدل غربالگری سرطان به‌جای فیلتر اسپم بود، دقت مثبت یا بازیابی را اولویت می‌دادید؟ چرا؟
\end{enumerate}

\bigskip

\textbf{(د)} [۴ نمره] یک طبقه‌بند احتمالات خروجی می‌دهد. آستانه طبقه‌بندی را از ۱.۰ تا ۰.۰ تغییر می‌دهید و منحنی ROC را رسم می‌کنید.

\begin{enumerate}[label=(\roman*)]
\item محور افقی و عمودی منحنی ROC چیست؟ فرمول آن‌ها را بنویسید.
\item تفسیر احتمالی AUC را در یک جمله بیان کنید.
\item اگر AUC $= 0.3$ به دست آمد، آیا این بدتر از حدس تصادفی است؟ چه باید بکنید؟
\end{enumerate}

\bigskip

\textbf{(هـ)} [۲ نمره] جای خالی را پر کنید: نمره F1 میانگین \underline{\hspace{3cm}} دقت مثبت و بازیابی است. در یک جمله توضیح دهید چرا این نوع میانگین به‌جای میانگین حسابی انتخاب شده است.

\newpage

\section*{بخش دوم: سوالات کوتاه (۳۰ دقیقه)}

\subsection*{سوال کوتاه ۱ [۶ نمره]}

$\mathcal{H}_1$ مجموعه توابع ثابت $f(x) = b$ و $\mathcal{H}_2$ مجموعه توابع خطی $f(x) = wx + b$ باشد. توجه کنید $\mathcal{H}_1 \subset \mathcal{H}_2$. فرض کنید $\theta_1^*$ و $\theta_2^*$ کمینه‌کننده‌های ریسک تجربی در هر فضا باشند.

\textbf{(الف)} ثابت کنید $R_{\text{emp}}(\theta_2^*) \leq R_{\text{emp}}(\theta_1^*)$.

\textbf{(ب)} آیا این بدان معناست که $\theta_2^*$ تعمیم بهتری دارد (ریسک واقعی کمتری)؟ توضیح دهید چرا.

\bigskip

\subsection*{سوال کوتاه ۲ [۶ نمره]}

برای جواب OLS، دو ویژگی کلیدی باقیمانده‌ها $e_i = y_i - \hat{y}_i$:

\begin{enumerate}
\item $\sum_{i=1}^n e_i = 0$ (مجموع باقیمانده‌ها صفر است)
\item $\sum_{i=1}^n x_i \cdot e_i = 0$ (باقیمانده‌ها با $x$ همبسته نیستند)
\end{enumerate}

\textbf{(الف)} به‌طور خلاصه توضیح دهید چرا هر ویژگی برقرار است.

\textbf{(ب)} با استفاده از ویژگی ۲، نشان دهید $\text{Cov}(\hat{y}, e) = 0$.

\textbf{(ج)} با استفاده از (ب)، تجزیه واریانس را ثابت کنید: $\text{Var}(y) = \text{Var}(\hat{y}) + \text{Var}(e)$.

\bigskip

\subsection*{سوال کوتاه ۳ [۶ نمره]}

هر بخش را در ۲-۳ جمله پاسخ دهید:

\textbf{(الف)} چرا لاسو جواب تحلیلی (closed-form) ندارد ولی ریج دارد؟ به مشتق‌پذیری ارجاع دهید.

\textbf{(ب)} شبکه الاستیک (elastic net) چیست و چه مشکلی از لاسوی خالص را حل می‌کند؟

\textbf{(ج)} عرض از مبدأ $b$ معمولاً در ریج یا لاسو تنظیم نمی‌شود. یک دلیل بیان کنید.

\bigskip

\subsection*{سوال کوتاه ۴ [۶ نمره]}

۵۰ نقطه داده دارید و باید بین اعتبارسنجی متقاطع ۵-بخشی و اعتبارسنجی Leave-One-Out (LOOCV) یکی را انتخاب کنید.

\textbf{(الف)} هر روش چند بار برازش مدل نیاز دارد؟

\textbf{(ب)} LOOCV بایاس کمتر ولی واریانس بالاتری نسبت به CV ۵-بخشی دارد. توضیح دهید چرا واریانس بالاتر است.

\textbf{(ج)} قانون طلایی ارزیابی مدل را بیان کنید. یک مثال از نشت داده و نحوه جلوگیری از آن بزنید.

\textbf{(د)} داده‌ها را با محاسبه میانگین و انحراف معیار روی همه داده‌ها نرمال‌سازی می‌کنید، سپس به آموزش/تست تقسیم می‌کنید. آیا این مشکل دارد؟ اگر بله، چگونه اصلاح می‌کنید؟

\bigskip

\subsection*{سوال کوتاه ۵ [۶ نمره]}

منحنی یادگیری زیر (خطای آموزش و خطای اعتبارسنجی بر حسب اندازه مجموعه آموزش) را در نظر بگیرید:

\begin{center}
\begin{tikzpicture}[scale=0.8]
\draw[->] (0,0) -- (7,0) node[right] {اندازه مجموعه آموزش};
\draw[->] (0,0) -- (0,4.5) node[above] {خطا};
\draw[dashed] (0,3.5) .. controls (3,3.2) and (4,2.8) .. (6.5,2.6) node[right] {\small اعتبارسنجی};
\draw[dashed] (0,1) .. controls (3,1.8) and (4,2.2) .. (6.5,2.4) node[right] {\small آموزش};
\end{tikzpicture}
\end{center}

\textbf{(الف)} تشخیص: کم‌برازش، بیش‌برازش یا برازش خوب؟ توجیه کنید.

\textbf{(ب)} آیا افزودن داده آموزش بیشتر کمک قابل‌توجهی می‌کند؟ چرا؟

\textbf{(ج)} افزایش یا کاهش $\lambda$ کمک می‌کند؟ توضیح دهید.

\textbf{(د)} این منحنی را با منحنی پیچیدگی (هفته ۳) مقایسه کنید: هر یک چه چیزی را تغییر می‌دهد و چگونه مکمل هم هستند؟

\newpage

\section*{پاسخ‌نامه}

\subsection*{پاسخ سوال ۱}

\textbf{(الف)} ریسک تجربی:
$$R_{\text{emp}}(w, b) = \frac{1}{n}\sum_{i=1}^{n}(wx_i + b - y_i)^2$$

ریسک واقعی:
$$R(w, b) = \mathbb{E}_{(x,y) \sim \mathcal{P}}\left[(wx + b - y)^2\right]$$

\textbf{اصل ERM:} یافتن پارامترهایی که میانگین زیان روی داده‌های آموزش را کمینه کنند: $\theta^* = \arg\min_{\theta} R_{\text{emp}}(\theta)$.

\textbf{چرا ریسک واقعی قابل کمینه‌کردن نیست:} ریسک واقعی به توزیع مشترک ناشناخته $\mathcal{P}$ روی $(x, y)$ بستگی دارد. ما فقط نمونه محدود $\mathcal{D}$ داریم و نمی‌توانیم امید ریاضی را محاسبه کنیم. ریسک تجربی تقریبی از ریسک واقعی است.

\bigskip

\textbf{(ب)} $w$ را ثابت فرض کنید. تابع زیان بر حسب $b$:
$$R(b) = \frac{1}{n}\sum_{i=1}^n (wx_i + b - y_i)^2$$

فرض کنید $\epsilon_i = wx_i - y_i$:
$$R(b) = \frac{1}{n}\sum_i (\epsilon_i + b)^2 = \frac{1}{n}\sum_i \epsilon_i^2 + 2b \cdot \frac{1}{n}\sum_i \epsilon_i + b^2$$

این یک درجه دوم بر حسب $b$ است: $R(b) = b^2 + 2b\bar{\epsilon} + \overline{\epsilon^2}$ که در آن $\bar{\epsilon} = \frac{1}{n}\sum_i \epsilon_i$.

درجه دوم $f(b) = b^2 + cb + d$ در $b = -c/2$ کمینه می‌شود، پس:
$$b^* = -\bar{\epsilon} = -\frac{1}{n}\sum_i (wx_i - y_i) = \bar{y} - w\bar{x}$$

\textbf{نقطه میانگین‌ها:} چون $b^* = \bar{y} - w\bar{x}$، داریم $\hat{y}(\bar{x}) = w\bar{x} + b^* = \bar{y}$. یعنی پیش‌بینی در $x = \bar{x}$ برابر $\bar{y}$ است، پس خط از $(\bar{x}, \bar{y})$ عبور می‌کند.

\bigskip

\textbf{(ج)} با جایگذاری $b^* = \bar{y} - w\bar{x}$، پیش‌بینی می‌شود $\hat{y}_i = w(x_i - \bar{x}) + \bar{y} = w\tilde{x}_i + \bar{y}$ که $\tilde{x}_i = x_i - \bar{x}$.

خطا: $\hat{y}_i - y_i = w\tilde{x}_i - \tilde{y}_i$ که $\tilde{y}_i = y_i - \bar{y}$.

$$R(w) = \frac{1}{n}\sum_i (w\tilde{x}_i - \tilde{y}_i)^2 = w^2 \underbrace{\frac{1}{n}\sum_i \tilde{x}_i^2}_{\text{Var}(x)} - 2w \underbrace{\frac{1}{n}\sum_i \tilde{x}_i\tilde{y}_i}_{\text{Cov}(x,y)} + \underbrace{\frac{1}{n}\sum_i \tilde{y}_i^2}_{\text{Var}(y)}$$

این درجه دوم $Aw^2 - Bw + C$ با $A = \text{Var}(x)$ و $B = 2\text{Cov}(x,y)$ است. کمینه در $w = B/(2A)$:
$$w^* = \frac{\text{Cov}(x,y)}{\text{Var}(x)}$$

\bigskip

\textbf{(د)}
$$w^*_{\text{ridge}} = \frac{\text{Cov}(x,y)}{\text{Var}(x) + \lambda} = w^*_{\text{OLS}} \cdot \underbrace{\frac{\text{Var}(x)}{\text{Var}(x) + \lambda}}_{\text{ضریب جمع‌شونده}}$$

ضریب جمع‌شونده $s = \frac{\text{Var}(x)}{\text{Var}(x) + \lambda} \in (0,1)$ برای $\lambda > 0$.

\begin{itemize}
\item $\lambda \to 0$: $s \to 1$، یعنی $w^*_{\text{ridge}} \to w^*_{\text{OLS}}$.
\item $\lambda \to \infty$: $s \to 0$، یعنی $w^*_{\text{ridge}} \to 0$ (مدل $\bar{y}$ را پیش‌بینی می‌کند).
\end{itemize}

\bigskip

\textbf{(هـ)}

\begin{center}
\begin{tabular}{|l|c|c|c|}
\hline
 & $\lambda = 0$ (OLS) & $\lambda$ متوسط & $\lambda \to \infty$ \\
\hline
\textbf{بایاس} & پایین & متوسط & بالا \\
\hline
\textbf{واریانس} & بالا & متوسط & پایین \\
\hline
\textbf{خطای آموزش} & کمترین & کمی بالاتر & بالاتر ($\approx \text{Var}(y)$) \\
\hline
\textbf{خطای تست} & ممکن است بالا & کمترین & بالا \\
\hline
\end{tabular}
\end{center}

\textbf{توضیح:} OLS بایاس پایین (برازش تهاجمی داده) ولی واریانس بالا (برازش نوفه، ناپایداری با داده‌های مختلف) دارد. ریج بایاس کمی را وارد می‌کند (جمع‌شدن به سمت صفر) اما واریانس را به‌طور چشمگیری کاهش می‌دهد (پایداری بیشتر). در نقطه بهینه، کاهش واریانس از افزایش بایاس بیشتر است، که به خطای تست کمتر منجر می‌شود.

\newpage

\subsection*{پاسخ سوال ۲}

\textbf{(الف)}
\begin{itemize}
\item \textbf{درجه ۱:} هر دو خطا بالا (حدود ۱۲-۱۳) و شکاف کوچک (۰.۶). \textbf{کم‌برازش} — مدل بیش از حد ساده است.
\item \textbf{درجه ۲:} هر دو خطا پایین (حدود ۱.۸-۲.۳) و شکاف کوچک (۰.۵). \textbf{برازش خوب} — الگوی درجه ۲ واقعی را ثبت می‌کند.
\item \textbf{درجه ۱۵:} خطای آموزش $\approx 0$ ولی خطای تست بسیار بالا (۳۲.۰) و شکاف عظیم. \textbf{بیش‌برازش} — نوفه را حفظ کرده و بد تعمیم می‌یابد.
\end{itemize}

\bigskip

\textbf{(ب)} چندجمله‌ای درجه $d_1$ حالت خاصی از چندجمله‌ای درجه $d_2$ (برای $d_1 < d_2$) است: کافی است ضرایب اضافی را صفر قرار دهیم. پس $\mathcal{H}_{d_1} \subset \mathcal{H}_{d_2}$.

هنگام کمینه‌سازی ریسک تجربی روی $\mathcal{H}_{d_2}$، بهینه‌ساز همواره می‌تواند جواب درجه $d_1$ را انتخاب کند. پس:
$$\min_{f \in \mathcal{H}_{d_2}} R_{\text{emp}}(f) \leq \min_{f \in \mathcal{H}_{d_1}} R_{\text{emp}}(f)$$

\bigskip

\textbf{(ج)} هر بخش: $20/5 = 4$ نمونه در اعتبارسنجی، $16$ نمونه در آموزش. مجموع برازش‌ها: ۵ (یک بار برای هر بخش) $\times$ تعداد درجه‌های کاندید. اگر ۶ درجه مقایسه شود: $5 \times 6 = 30$ برازش.

\textbf{مراحل ۵گانه:}
\begin{enumerate}
\item انتخاب درجه‌های کاندید (مثلاً $d = 1, 2, 3, 5, 8, 15$).
\item برای هر درجه $d$، اعتبارسنجی متقاطع ۵-بخشی: داده آموزش به ۵ بخش، آموزش روی ۴ بخش، ارزیابی روی ۱ بخش، تکرار ۵ بار، میانگین خطاها.
\item انتخاب درجه با کمترین میانگین خطای CV.
\item (اختیاری) بازآموزش روی همه داده آموزش با درجه انتخاب‌شده.
\item ارزیابی \textbf{یک‌بار} روی مجموعه تست.
\end{enumerate}

\bigskip

\textbf{(د)} \textbf{بیش‌برازش به مجموعه اعتبارسنجی:} وقتی مدل‌های زیادی (۵۰۰) امتحان می‌شوند و بهترین روی مجموعه اعتبارسنجی انتخاب می‌شود، برنده تا حدی به‌خاطر شانس خوب است — مدلی به‌طور تصادفی روی اعتبارسنجی خوب دیده می‌شود. این مشکل مقایسات چندگانه است. خطای اعتبارسنجی برنده به‌طور خوش‌بینانه‌ای بایاس‌دار است.

\textbf{چرا تست کمک می‌کند:} مجموعه تست بعد از تمام انتخاب‌ها، فقط یک‌بار استفاده می‌شود. چون مدلی بر اساس آن انتخاب نشده، بایاس انتخابی وجود ندارد و تخمین صادقانه‌ای از عملکرد تعمیم می‌دهد.

\bigskip

\textbf{(هـ)} \textbf{نشت زمانی:} مخلوط کردن داده‌های سری زمانی نظم زمانی را می‌شکند، پس مجموعه آموزش ممکن است داده‌های آینده نسبت به تست را داشته باشد. در استقرار واقعی نمی‌توان آینده را دید، پس این تخمین بیش از حد خوش‌بینانه می‌دهد.

\textbf{رویکرد صحیح:} آموزش روی گذشته، تست روی آینده. برای CV از سری زمانی: آموزش روی $[1..t]$، اعتبارسنجی روی $[t+1..t+k]$، سپس بسط پنجره آموزش.

\bigskip

\textbf{(و)} $R^2 = 0.91$ یعنی مدل ۹۱٪ واریانس $y$ را توضیح می‌دهد: $R^2 = 1 - \text{SS}_{\text{res}}/\text{SS}_{\text{tot}}$، پس جمع مربعات باقیمانده فقط ۹٪ جمع مربعات کل است.

$R^2 = -0.05$ یعنی مدل \textbf{بدتر} از پیش‌بینی $\bar{y}$ برای همه ورودی‌ها عمل می‌کند. جمع مربعات باقیمانده از جمع مربعات کل بیشتر است. این می‌تواند در اثر بیش‌برازش روی داده آموزش رخ دهد.

\newpage

\subsection*{پاسخ سوال ۳}

\textbf{(الف)} \textbf{(i)} خطای مربعی ($L_2$) حساس‌تر است. مثال: خطاها $\{1, 1, 1, 1, 10\}$.
\begin{itemize}
\item MSE: سهم داده پرت $= 10^2 = 100$، کل $= (1+1+1+1+100)/5 = 20.8$. داده پرت غالب است.
\item MAE: سهم داده پرت $= 10$، کل $= (1+1+1+1+10)/5 = 2.8$. متناسب، نه غالب.
\end{itemize}
مربع کردن خطاهای بزرگ را تشدید می‌کند.

\textbf{(ii)} هوبر برای خطاهای کوچک از $L_2$ (هموار، مشتق‌پذیر، بهینه‌سازی آسان) و برای خطاهای بزرگ از $L_1$ (مقاوم در برابر داده پرت) استفاده می‌کند. پیش‌بینی‌های معمولی از خواص خوب $L_2$ بهره می‌برند و داده‌های پرت مانند $L_2$ خالص غالب نمی‌شوند.

\bigskip

\textbf{(ب)} \textbf{(i)}
\begin{itemize}
\item ریج (L2): $\lambda(w_1^2 + w_2^2)$
\item لاسو (L1): $\lambda(|w_1| + |w_2|)$
\end{itemize}

\textbf{(ii)}
\begin{itemize}
\item محدودیت L2 ($\|\mathbf{w}\|^2 \leq t$): یک \textbf{دایره} به مرکز مبدأ.
\item محدودیت L1 ($\|\mathbf{w}\|_1 \leq t$): یک \textbf{لوزی} با رئوس روی محورها در $(\pm t, 0)$ و $(0, \pm t)$.
\end{itemize}

\textbf{(iii)} جواب OLS در فضای داخلی است. تنظیم آن را به نزدیک‌ترین نقطه روی ناحیه محدودیت می‌کشد. برای دایره L2، نزدیک‌ترین نقطه معمولاً \textbf{روی محور نیست}، پس هر دو $w_1$ و $w_2$ ناصفر (ولی کوچک) می‌مانند. برای لوزی L1، نزدیک‌ترین نقطه اغلب یک \textbf{رأس} است — و رئوس روی محورها هستند، جایی که یکی از $w_j = 0$ است. این دلیل تنکی لاسو است.

\textbf{(iv)}
\begin{itemize}
\item لاسو: وقتی ویژگی‌های زیادی دارید ولی فقط چند مورد مؤثر است (انتخاب ویژگی).
\item ریج: وقتی همه ویژگی‌ها مؤثرند یا ویژگی‌ها همبسته‌اند.
\end{itemize}

\bigskip

\textbf{(ج)} $TP = 160$، $FP = 30$، $FN = 40$، $TN = 1770$. کل $= 2000$.

\textbf{(i)}
\begin{itemize}
\item \textbf{دقت} $= (160 + 1770)/2000 = 1930/2000 = 96.5\%$
\item \textbf{دقت مثبت} $= 160/(160+30) = 160/190 = 84.2\%$
\item \textbf{بازیابی} $= 160/(160+40) = 160/200 = 80.0\%$
\item \textbf{F1} $= 2 \times (0.842 \times 0.80)/(0.842 + 0.80) = 82.0\%$
\end{itemize}

\textbf{(ii)} مدل پایه (همه غیراسپم): $TP = 0$، دقت $= 1800/2000 = 90\%$. پایه با هیچ‌کاری ۹۰٪ دقت می‌گیرد. مدل ۹۶.۵٪ به نظر خوب می‌آید، ولی پایه ۹۰٪ است. بهبود واقعی در بازیابی (۸۰٪ در برابر ۰٪) و F1 است. دقت توسط کلاس اکثریت غالب می‌شود.

\textbf{(iii)} برای غربالگری سرطان، \textbf{بازیابی} اولویت دارد. از دست دادن مورد سرطان (منفی کاذب) بالقوه کشنده است، در حالی که هشدار کاذب (مثبت کاذب) فقط آزمایش اضافی و اضطراب دارد.

\bigskip

\textbf{(د)} \textbf{(i)}
\begin{itemize}
\item محور افقی: FPR $= \frac{FP}{FP + TN}$
\item محور عمودی: TPR $= \frac{TP}{TP + FN}$ (همان بازیابی)
\end{itemize}

\textbf{(ii)} AUC = احتمال اینکه طبقه‌بند، یک مثبت تصادفی را بالاتر از یک منفی تصادفی رتبه‌بندی کند.

\textbf{(iii)} AUC $= 0.3$ \textbf{بدتر} از حدس تصادفی ($0.5$) است. مدل به‌طور سیستماتیک منفی‌ها را بالاتر از مثبت‌ها رتبه‌بندی می‌کند. باید \textbf{پیش‌بینی‌ها را معکوس} کنید که AUC $= 0.7$ می‌دهد.

\bigskip

\textbf{(هـ)} میانگین \textbf{هارمونیک}.

میانگین هارمونیک عدم تعادل شدید را تنبیه می‌کند: اگر دقت مثبت $= 0.01$ و بازیابی $= 1.0$، میانگین حسابی $= 0.505$ (خوب به نظر می‌رسد) ولی میانگین هارمونیک $\approx 0.02$ (درست، افتضاح). میانگین هارمونیک فقط وقتی بالا است که \textbf{هر دو} بالا باشند.

\newpage

\subsection*{پاسخ سوالات کوتاه}

\subsubsection*{پاسخ سوال کوتاه ۱}

\textbf{(الف)} چون $\mathcal{H}_1 \subset \mathcal{H}_2$، هر تابع در $\mathcal{H}_1$ در $\mathcal{H}_2$ نیز هست. کمینه‌سازی روی $\mathcal{H}_2$ جستجو در یک \textbf{ابرمجموعه} از $\mathcal{H}_1$ است. کمینه روی مجموعه بزرگتر حداکثر برابر کمینه روی مجموعه کوچکتر است:
$$R_{\text{emp}}(\theta_2^*) \leq R_{\text{emp}}(\theta_1^*)$$

\textbf{(ب)} خیر. ریسک تجربی کمتر، تضمینی برای ریسک واقعی کمتر نیست. مدل پیچیده‌تر $\theta_2^*$ ممکن است نوفه داده آموزش را برازش داده باشد (بیش‌برازش)، که به شکاف تعمیم بزرگ منجر می‌شود. این دقیقاً مسئله بیش‌برازش در برابر کم‌برازش است.

\bigskip

\subsubsection*{پاسخ سوال کوتاه ۲}

\textbf{(الف)}
\begin{enumerate}
\item $\sum_i e_i = 0$: از فرمول عرض از مبدأ $b^* = \bar{y} - w^*\bar{x}$ نتیجه می‌شود. خط طوری قرار می‌گیرد که خطاهای مثبت و منفی متعادل شوند.
\item $\sum_i x_i e_i = 0$: جواب OLS \textbf{تمام} رابطه خطی بین $x$ و $y$ را ثبت می‌کند. باقیمانده‌ها هیچ الگوی خطی در $x$ ندارند.
\end{enumerate}

\textbf{(ب)} چون $\hat{y}_i = w^*x_i + b^*$ تابعی خطی از $x_i$ است:
$$\sum_i \hat{y}_i e_i = w^* \sum_i x_i e_i + b^* \sum_i e_i = 0 + 0 = 0$$
پس $\text{Cov}(\hat{y}, e) = 0$.

\textbf{(ج)} چون $y_i = \hat{y}_i + e_i$:
$$\text{Var}(y) = \text{Var}(\hat{y} + e) = \text{Var}(\hat{y}) + \text{Var}(e) + 2\text{Cov}(\hat{y}, e) = \text{Var}(\hat{y}) + \text{Var}(e)$$

\bigskip

\subsubsection*{پاسخ سوال کوتاه ۳}

\textbf{(الف)} جریمه L1 یعنی $|w|$ در $w = 0$ مشتق‌پذیر نیست (مشتق چپ $-1$، مشتق راست $+1$). نمی‌توان گرادیان را صفر قرار داد. اما $w^2$ ریج هموار و مشتق‌پذیر است.

\textbf{(ب)} شبکه الاستیک ترکیب L1 و L2: $\lambda_1\sum|w_j| + \lambda_2\sum w_j^2$. مشکل لاسو با ویژگی‌های همبسته را حل می‌کند: لاسو یکی از دو ویژگی همبسته را تصادفی انتخاب و دیگری را صفر می‌کند، اما شبکه الاستیک هر دو را نگه می‌دارد.

\textbf{(ج)} عرض از مبدأ $b$ فقط جابجایی عمودی خط است، نه حساسیت به $x$. بیش‌برازش از شیب‌های بزرگ ناشی می‌شود، نه از عرض از مبدأ.

\bigskip

\subsubsection*{پاسخ سوال کوتاه ۴}

\textbf{(الف)} CV ۵-بخشی: ۵ برازش. LOOCV: ۵۰ برازش.

\textbf{(ب)} در LOOCV هر مجموعه آموزش $n-1 = 49$ نقطه دارد — مجموعه‌های آموزش تقریباً یکسان هستند. ۵۰ خطای اعتبارسنجی به‌شدت همبسته‌اند. میانگین کمیت‌های همبسته واریانس را چندان کاهش نمی‌دهد.

\textbf{(ج)} \textbf{قانون طلایی:} هرگز در طول توسعه مدل به داده تست دست نزنید. مجموعه تست فقط یک‌بار برای ارزیابی نهایی استفاده می‌شود.

مثال نشت: نرمال‌سازی با میانگین و انحراف معیار \textbf{همه} داده‌ها قبل از تقسیم. \textbf{راه‌حل:} اول تقسیم کنید، سپس میانگین و انحراف معیار را فقط روی داده آموزش محاسبه و به داده تست اعمال کنید.

\textbf{(د)} بله، این نشت داده است. میانگین و انحراف معیار شامل داده تست است. \textbf{اصلاح:} اول تقسیم کنید، سپس نرمال‌سازی را فقط با آمار داده آموزش انجام دهید.

\bigskip

\subsubsection*{پاسخ سوال کوتاه ۵}

\textbf{(الف)} \textbf{برازش خوب} (یا کم‌برازش بسیار خفیف). هر دو خطا پایین و شکاف کوچک و در حال کاهش است. منحنی‌ها تقریباً همگرا شده‌اند.

\textbf{(ب)} کمک حاشیه‌ای. منحنی‌ها همگرا شده‌اند، یعنی مدل بیشتر سیگنال قابل‌یادگیری را استخراج کرده. خطای باقیمانده احتمالاً غیرقابل‌کاهش (نوفه) یا ناشی از بایاس مدل است. داده بیشتر عمدتاً واریانس را کاهش می‌دهد، ولی شکاف از قبل کوچک است.

\textbf{(ج)} چون مدل برازش خوبی با شکاف کوچک دارد، $\lambda$ احتمالاً خوب تنظیم شده. اگر کم‌برازش خفیفی وجود دارد، \textbf{کاهش} $\lambda$ (تنظیم کمتر) ممکن است کمکی کند، ولی اثر کم خواهد بود.

\textbf{(د)} \textbf{منحنی پیچیدگی} (هفته ۳): اندازه داده ثابت، پیچیدگی مدل تغییر می‌کند (درجه، $\lambda$). منحنی U شکل خطای تست را نشان می‌دهد و نقطه بهینه پیچیدگی را پیدا می‌کند. \textbf{منحنی یادگیری} (هفته ۴): مدل ثابت، اندازه داده تغییر می‌کند. نشان می‌دهد آیا داده بیشتر کمک می‌کند (شکاف بزرگ → بله) یا مدل به حد خود رسیده (همگرا → خیر). مکمل‌اند: منحنی پیچیدگی می‌گوید \textit{چه مدلی}، منحنی یادگیری می‌گوید \textit{آیا داده بیشتر}.

\end{document}
